
\svnid{$Id: case_text.tex 24504 2019-08-27 16:03:25Z hutten $}
\svnid{$URL: https://repos.deltares.nl/repos/DSCTestbench/trunk/cases/e02_dflowfm/f105_cross-sections/c16_YZ-profile_lumped/doc/chapters/case_text.tex $}

\section{YZ-profile using lumped conveyance approach}
\newrefsegment

%---------------------------------------------------------------------------------
\section*{Quality Assurance}
\begin{tabular}{@{}p{27mm}@{}p{0mm}p{\textwidth-52mm-48pt}p{15mm}p{10mm}}
\textbf{Date} & \textbf{:} &   July 22, 2019\\
\textbf{Author} & \textbf{:} & Christianne Luger & \textbf{Initials: } & CL \\
\textbf{Review} & \textbf{:} &  Bert Jagers & \textbf{Initials:} & BJ 
\end{tabular}

%---------------------------------------------------------------------------------
\section*{Version information}
{{
\begin{tabular}{@{}p{27mm}@{}p{0mm}p{\textwidth-27mm-24pt}}
\textbf{Executable}   & \textbf{:} & Deltares, D-Flow FM Version 1.2.59.64502M, Aug 02 2019, 10:44:19  \\
\textbf{Location input}     & \textbf{:} & \url{https://repos.deltares.nl/repos/DSCTestbench/trunk/cases/e02_dflowfm/f105_cross-sections/c16_YZ-profile_lumped}  \\
\textbf{Revision input} & \textbf{:} & - \\
\textbf{Location output}     & \textbf{:} & \url{https://repos.deltares.nl/repos/DSCTestbench/trunk/references/win64/e02_dflowfm/f105_cross-sections/c16_YZ-profile_lumped/DFM_OUTPUT_Lumped_YZ_flow_area} \\
\textbf{Revision output} & \textbf{:} & -
\end{tabular}
}}

%---------------------------------------------------------------------------------
\section*{Purpose}
The purpose of this validation test is to verify that \dflowfm gives the correct result with the implementation of YZ-profiles using the lumped conveyance approach and that upstream water depth shows right balance between roughness and resistance while applying different roughness types.

%---------------------------------------------------------------------------------
\section*{Linked claims}
\begin{enumerate}
\item \dflowfm can take into account the shape of several cross sectional profiles (circle, rectangle, ZW, XYZ and YZ).
\end{enumerate}

%---------------------------------------------------------------------------------
\section*{Approach}

YZ-profiles using lumped conveyance method are validated for different roughness types by verifying that the upstream water depth of the channels show the right balance between roughness and resistance. 
The model consist of 6 sub-models T1 to T6, each sub-model in \autoref{fig:e02_f105_c16_Test model lay-out} is parametrized with a different roughness value and type. 
For each sub-model, the upstream conveyance ($K(h)$) is analytically computed for a {constant} water depth ($h$) of 3 m using the equations for the lumped approach. 
This calculated conveyance is converted to a {constant} discharge ($Q$) condition at equilibrium using the equation $Q=K(h)\sqrt{S}$ where $S$ is the slope of the channel. 
To verify the claims for this case, the modeled water levels throughout the sub-models must be steady and equal to 3 m if the equilibrium discharge computed is routed through the systems.


%---------------------------------------------------------------------------------
\section*{Conclusion} 
The implementation of YZ-profiles using lumped conveyance method for steady state simulations while applying a downstream water level boundary close to equilibrium depth is considered to be satisfactory. 
This conclusion is valid for the application of the following roughness coefficients:
\begin{itemize}
\item Ch\'ezy
\item Manning
\item Strickler ks
\item Nikuradse (Strickler)
\item Nikuradse (White \& Colebrook)
\item De Bos \& Bijkerk
\end{itemize}

The water depths at all sub-models show the right balance between roughness and resistance. 
The differences between the modelled and analytical solutions are due to the shifts introduced by the calculations in the cross-sectional profile when there exists equal and adjacent $z$-values. 
These differences however are minimal and within the threshold of acceptance.

Because the downstream water level boundary condition cannot match the equilibrium conditions exactly, the model includes a backwater effect and the equilibrium condition is only reached approximately 10 km away from the boundary.

%---------------------------------------------------------------------------------
\section*{Model description}
The test comprises of six separate sub-models {T1} to {T6} (see \autoref{fig:e02_f105_c16_Test model lay-out}). 
The sub-models are 20 km in length and have an equidistant grid spacing of dx=50 m. 
The YZ-profile is constant throughout the length of the models (see \autoref{fig:e02_f105_c16_YZCrossSection}).
The bed levels decrease by 5 m from x=0 (East, upstream) to x=20000 (West, downstream) for a slope of 0.00025. 
Roughness types and values vary among the sub-models but it is constant across the cross-section of a sub-model since the lumped conveyance approach is used (See \autoref{e02_f105_c16_RoughnessYZ}).

\begin{figure}[H]
    \centering
    \includegraphics*[width=0.95\textwidth]{figures/e02_f105_c16_layout.png}
    \caption{Lay-out of the test model}
    \label{fig:e02_f105_c16_Test model lay-out}
\end{figure}

\begin{figure}[H]
    \centering
    \includegraphics*[width=0.95\textwidth]{figures/e02_f105_c16_cross_section.png}
    \caption{Cross-section definition of YZ-profile at x=20000 (West, downstream)}
    \label{fig:e02_f105_c16_YZCrossSection}
\end{figure}

{Note} that \dflowfm does not allow for horizontal and vertical lines in the YZ-profile. 
Vertical lines cannot be specified in the user interface and therefore a small offset has been added to the $y'$ for the vertical sections at the left and right bank.
When two neighbouring $z$ levels are equal, then the second point is automatically shifted up by 1.1mm by the kernel.

\begin{longtable}{|p{18mm}|p{50mm}|p{30.0mm}|}
\caption{Roughness types and values of  {T1} to {T6}.} 
\label{e02_f105_c16_RoughnessYZ} \\ 
\hline
\STRUT
      {Sub-model} 
	 &  {Roughness Type } 
	 &  {Roughness Value} \\ 
[1ex]\hline \endhead
\hline
\endfoot
%
\STRUT
{T1} & {Ch\'ezy} & {45} \\
\STRUT
{T2} & {Manning} & {0.05} \\	
\STRUT
{T3} & {Strickler ks} & {40} \\
\STRUT
{T4} & {Nikuradse (Strickler)} & {0.2} \\
\STRUT
{T5} & {Nikuradse (White \& Colebrook)} & {0.3} \\
\STRUT
{T6} & {De Bos \& Bijkerk} & {30} \\ [1ex]	
\end{longtable}

The boundary conditions of all sub-models are listed in \autoref{e02_f105_c16_Boundary conditions}. 
The upstream flow boundary conditions are parametrized to ensure that all sub-models are at equilibrium with a water depth of 3 m. The upstream boundary condition varies for each sub-model as their roughness types and values are different. 
To ensure that water depth are at 3 m at x=20000, the downstream boundary condition is set to $3-(S*\frac{dx}{2})$. 
This is because water level boundaries are in \dflowfm not applied at the last grid point of the network, but depending on the setting of  \texttt{izbndpos} some distance outside the system. 
For these simulations we have set \texttt{izbndpos = 1} which means that the water levels are imposed half a computational grid cell length outside the system.

Note that water level boundaries are applied by default in \dflowfm one whole computational grid cell length outside the system. 

\begin{longtable}{|p{18mm}|p{23mm}|p{35.0mm}|p{65mm}|}
\caption{Boundary conditions applied in sub-models {T1} to{T6}.} 
\label{e02_f105_c16_Boundary conditions} \\ 
\hline
\STRUT
      {Sub-model} 
	 &  {Chainage \unitbrackets{m}}  
	 &  {ID}  
	&  {Boundary condition}  \\ 
[1ex]\hline \endhead
\hline
\endfoot
%
\STRUT
{T1} &  {0.00}     & {T1\_C\_x0m}     & {Constant discharge of 121.33173 m\textsuperscript{3}s\textsuperscript{-1}}   		      \\
                                               & {20000.00}        & {T1\_C\_x20000m}    & {Constant water level of 2.99375 m AD}                		      \\

\STRUT
{T2} & {0.00}     & {T2\_n\_x0m}     & {Constant discharge of 61.49307 m\textsuperscript{3}s\textsuperscript{-1}}   	      \\
                                               & {20000.00}        & {T2\_n\_x20000m}    & {Constant water level of 2.99375 m AD}                           \\

\STRUT
{T3} & {0.00}     & {T3\_ks\_x0m}     & {Constant discharge of 122.98613 m\textsuperscript{3}s\textsuperscript{-1}}       \\
                                               & {20000.00}        & {T3\_ks\_x20000m}    & {Constant water level of 2.99375 m AD}                      \\

\STRUT										   
{T4} & {0.00}     & {T4\_kn\_x0m}   & {Constant discharge of 100.515067 m\textsuperscript{3}s\textsuperscript{-1}}   \\
                                               & {20000.00}        & {T4\_kn\_x20000m}       & {Constant water level of 2.99375 m AD}                      \\
	
\STRUT										   
{T5} & {0.00}     & {T5\_WAQUA\_x0m}     & {Constant discharge of 94.36045 m\textsuperscript{3}s\textsuperscript{-1}}   \\
                                               & {20000.00}        & {T5\_WAQUA\_x20000m}         & {Constant water level of 2.99375 m AD}                     \\

\STRUT											   
{T6} & {0.00}     & {T6\_gamma\_x0m}  & {Constant discharge of 133.03252 m\textsuperscript{3}s\textsuperscript{-1}}   \\
                                               & {20000.00}        & {T6\_gamma\_x20000m}      & {Constant water level of 2.99375 m AD}                     \\ [1ex]
                                                                                                                              
\end{longtable}	 


%---------------------------------------------------------------------------------
\section*{Results}
The modelled water depths along sub-models {T1} to {T6} are shown in \autoref{fig:e02_f105_c16_WaterDepthEntireLength}. 
The modelled upstream water depths (x=0) and their difference from the expected analytical solution are given in \autoref{tab:e02_f105_c16__results}. 

\begin{figure}[H]
    \centering
    \includegraphics*[width=0.95\textwidth]{figures/e02_f105_c16_Water_Depth_EntireLength.jpg}
    \caption{Water depths along sub-models {T1} to {T6}}
    \label{fig:e02_f105_c16_WaterDepthEntireLength}
\end{figure} 

\begin{longtable}{|p{18mm}|p{20mm}|p{20mm}|p{35mm}|}
\caption{Upstream (x=0) modelled and expected analytical water depth $h$.} 
\label{tab:e02_f105_c16__results} \\ 
\hline
\STRUT
      {Sub-model} 
	 &  {$h$\textsubscript{modelled}} \unitbrackets{m} 
	 &  {$h$\textsubscript{analytical}} \unitbrackets{m} 
	 &  {$\Delta h$\textsubscript{modelled-analytical}} \unitbrackets{m} \\
[1ex] \hline \endhead
\hline
\endfoot
%
\STRUT
{{T1}} & {3.00050375} & {3} & {0.00050375} \\	
\STRUT
{{T2}} & {3.00050719} & {3} & {0.00050719} \\	
\STRUT
{{T3}} & {3.00050765} & {3} & {0.00050765} \\
\STRUT
{{T4}} & {3.00050746} & {3} & {0.00050746} \\
\STRUT
{{T5}} & {3.00050848} & {3} & {0.00050848} \\	
\STRUT
{{T6}} & {3.00043889} & {3} & {0.00043889} \\ [1ex]	
\end{longtable}	 

%---------------------------------------------------------------------------------
\section*{Analysis of results}
All sub-models converge to an equilibrium water depth approximately 0.5 mm higher than the expected analytical value. 
This difference can be explained by the automatic shifting of equal and adjacent $z$-values by 0.0011 m in the YZ-profile. 
This effect is confirmed by means of reference testcase f105\_cross-sections/c17\_YZ-profile\_lumped\_V-shaped.

Another observation from the results is that the equilibrium state is not reached downstream. 
The approximation of the downstream boundary condition is inaccurate because it is based on the assumption that it is applied half a grid cell outside the system and is extrapolated using the system slope, whereas the model actually assumes a horizontal (zero slope) extension beyond the end of the channel. 

%---------------------------------------------------------------------------------
\section*{Lumped conveyance approach} \label{e02_f105_c16_description of Lumped conveyance approach}
In the lumped conveyance approach, it is assumed that there is an uniform flow velocity (\=U) in the cross-sectional profile for a given water level (see \autoref{fig:10101:Lumped approach}).
Hence, it is assumed that differences in roughness coefficients and flow velocities across the cross-sectional profile can be neglected.

\begin{figure}[H]
    \centering
    \includegraphics*[width=0.65\textwidth]{figures/e02_f105_c16_lumped_conveyance.png}
    \caption{Concept of the lumped conveyance approach}
    \label{fig:10101:Lumped approach}
\end{figure}

The lumped conveyance $K$ is computed as a function of water depth $h$:

\begin{equation}
K(h) = A(h)C(h)\sqrt{R(h)}
\end{equation}

where:
\begin{itemize}
\item $K(h)$ Lumped conveyance at $h$
\item $A(h)$ Cross-sectional area at $h$
\item $C(h)$ Ch\'ezy friction value at $h$
\item $P(h)$ Wetted perimeter at $h$
\item $R(h)$ Hydraulic radius at $h(R(h)=A(h)/P(h))$
\end{itemize}

Note that the wetted perimeter $P$ as function of $h$ would be discontinuous when a part of the profile were horizontal; therefore, we don't allow for that situation.
Roughness types other than Ch\'ezy will be converted to their Ch\'ezy equivalent based on the local water depth $h$ and hydraulic radius $R(h)$.
%---------------------------------------------------------------------------------
\printrefsegment
